\subsubsection{A KKL-KMS (Kahn--Kalai--Linial --- Keller--Mossel--Sen) theorem for non-product Gaussian vectors}\label{sss.KMS}

One of the contributions of this paper is the derivation of a KKL theorem for non-product Gaussian vectors, namely Theorem~\ref{t.KMSnonproduct}, based on a similar result for product Gaussian vectors by Keller, Mossel and Sen~\cite{keller2012geometric}. (This similar result by~\cite{keller2012geometric} is stated in Theorem~\ref{t.KMS} of our paper\footnote{As we will explain, Theorem~\ref{t.KMS} is a simple consequence of Item~$(1)$ of Theorem~1.5 of~\cite{keller2012geometric}.}.) Actually, with our techniques of Section~\ref{s.KMS}, most of the results from~\cite{keller2012geometric} could be extended to non-product Gaussian vectors (and to monotonic events).

\begin{theo}\label{t.KMSnonproduct}
There exists an absolute constant $c > 0$ such that the following holds: Let $n \in \Z_{>0}$, let $\Sigma$ be a $n \times n$ symmetric positive definite matrix\footnote{Remember Remark~\ref{r.pos_def}: this means in particular that $\Sigma$ is non-degenerate.} and let $\mu = \mathcal{N}(0,\Sigma)$. Also, let $\sqrt{\Sigma}$ be a symmetric square root of $\Sigma$ and write $\| \sqrt{\Sigma} \|_{\infty,op}$ for the operator norm of $\sqrt{\Sigma}$ for the infinite norm\footnote{I.e. $\| \sqrt{\Sigma} \|_{\infty,op} \Mk=\Mk \sup_{x\in \R^n \setminus \{0\}} \frac{|\sqrt{\Sigma}(x) |_\infty}{| x |_\infty}$. Equivalently, $\| \sqrt{\Sigma} \|_{\infty,op} \Mk=\Mk\max_{i \in \lbrace 1,\dots, n \rbrace} \sum_{j=1}^n |\sqrt{\Sigma}(i,j)|$.} on $\R^n$. For every $A$ monotonic Borel subset of $\R^n$ we have:
\[
\sum_{i=1}^n I_{i,\mu}(A) \geq c \, \| \sqrt{\Sigma} \|_{\infty,op}^{-1} \, \mu(A) \, (1-\mu(A)) \, \sqrt{\log_+ \left(\frac{1}{\| \sqrt{\Sigma} \|_{\infty,op} \cdot \max_{i \in \lbrace 1,\dots, n \rbrace} I_{i,\mu}(A)} \right)} \,.
\]
\end{theo}

This theorem holds for a wider class of sets. For instance, it holds for semi-algebraic\footnote{We say that a set $A\subset\R^n$ is semi-algebraic if it belongs to the Boolean algebra generated by sets of the form $P^{-1}\left(]0,+\infty[\right)$ where $P\in\R[X_1,\dots,X_n]$.}, see Appendix A of~\cite[Chapter~7]{alelele} or Appendix A of~\cite[Chapter~2]{hugogogo}. Since we need it only for monotonic events, we did not try to identify the weakest possible assumptions for the property to hold. In particular, we have not found any examples of Borel sets for which the theorem does not hold.

A way to understand the constant $\| \sqrt{\Sigma} \|_{\infty,op}$ is to see what happens in the extreme cases:
\begin{itemize}
\item If $\Sigma$ is the identity matrix, then $\| \sqrt{\Sigma} \|_{\infty,op}=1$ and the above result corresponds to the product case from~\cite{keller2012geometric}.
\item If $\Sigma_{i,j}=1$ for every $i,j$ (which corresponds to the fact that, if $X \sim \Sigma$, then $X_i=X_j$ for every $i,j$) then $\| \sqrt{\Sigma} \|_{\infty,op}=n$ which is coherent since there is no threshold phenomenon for a single variable.
\end{itemize}

The proof of Theorem~\ref{t.KMSnonproduct} is organized as follows: In Subsection~\ref{ss.KMSnonproduct}, we explain how to deduce Theorem~\ref{t.KMSnonproduct} from a sub-linear property of the influences (see Proposition~\ref{p.sublin}) and from the results of~\cite{keller2012geometric}. In Subsection~\ref{ss.monotonic}, we prove the sub-linearity property for monotonic subsets.
